\section{TODO}
\begin{enumerate}
    \item \sout{Read \cite{farber2006topology}}
    \item \sout{Read \cite{colman2013equivariant}}
    \item Read \cite{dym2024equivariant}
    \item \sout{Ok so there's no continuous section in \cite{dym2024equivariant}, but what is the Schwarz genus? I.e. how many cases are there?}
    \item \sout{Is it a good idea to just canonicalize via $\S^n_d$? Why or why not?}
    \item \sout{What can we learn about $X = \set{A \in \R^{d \times n} : \rank(A) = d}$ and $X/O(d)$ or $X/SO(d)$?}
    \item Is a canonicalization up to homotopy sufficient?
    \item \sout{Use restrictions to bound sectional category on original space}
    \item \sout{Consider $\SO(d)$ acting on $\R^{d \times n}$}
    \item \sout{Understand proof that there is no canonicalization of $(\R^{d \times n}, \O(d))$ and $(\R^{d \times n}, \SO(d))$}.
    \item \sout{What if we want a \emph{differentiable} section?}
    \item \sout{Read up on frames}
    \item Read about equivariant neural networks/ML to provide advantages/disadvantages and motivation for canonicalization/frames
    \item Formalize Richter-Tchakaloff approach to creating canonicalization frames
        \subitem Generalize? I.e. we don't actually need to compute $\int_G f(\inv{g}x)dg$, just some invariant map.
    \item Apply \cite{dym2024equivariant}'s results on frames to \cite{ma2024canonicalizationperspectiveinvariantequivariant}'s canonicalization frames
    \item Consider probability measures on a (finite) group $G$ as convex combinations of the simplex on $G$, and see what can be learned about weighted frames
    \item What can we say about the cardinality of continuity-preserving frames for $\O(d)$ and $\SO(d)$ for $d > 2$?
    \item Find (canonicalization) frames and/or bounds for (canonicalization) frame cardinality for more groups/spaces
    \item Construct explicit piecewise continuous sections for:
        \subitem $(\R^{d \times n},\O(d))$
        \subitem $(\R^{d \times n},\SO(d))$
        \subitem $(\R^{d \times n},\Sigma_n)$
        \subitem $(\R^{d \times n}/\Sigma_n,\O(d))$
        \subitem $(\R^{d \times n}/\Sigma_n,\SO(d))$
    \item Experiment computationally with piecewise continuous sections
    \item Consider the case of rotations acting on images
\end{enumerate}